\Chapter{87}

\Verse{1}
{
    \zR{却说黛玉叫进宝钗家的女人来，问了好，呈上书子，黛玉叫他去喝茶，便将宝 钗来书打开看时，只见上面写着： 妹生辰不偶，家运多艰，姊妹伶仃，萱亲衰迈。}
}
{
    \eR{[English source not present in the checked-in Bencraft/Joly files.]}
}
{
}

\Verse{2}
{
    \zR{兼之？}
}
{
    \eR{[English source not present in the checked-in Bencraft/Joly files.]}
}
{
}

\Verse{3}
{
    \zR{声狺语，旦暮无休；更 遭惨祸飞灾，不啻惊风密雨。}
}
{
    \eR{[English source not present in the checked-in Bencraft/Joly files.]}
}
{
}

\Verse{4}
{
    \zR{夜深辗侧，愁绪何堪。}
}
{
    \eR{[English source not present in the checked-in Bencraft/Joly files.]}
}
{
}

\Verse{5}
{
    \zR{属在同心，能不为之愍恻乎？}
}
{
    \eR{[English source not present in the checked-in Bencraft/Joly files.]}
}
{
}

\Verse{6}
{
    \zR{回忆海棠结社，序属清秋，对菊持螯，同盟欢洽。}
}
{
    \eR{[English source not present in the checked-in Bencraft/Joly files.]}
}
{
}

\Verse{7}
{
    \zR{犹记“孤标傲世偕谁隐，一样开 花为底迟”之句，未尝不叹冷节馀芳，如吾两人也!感怀触绪，聊赋四章。}
}
{
    \eR{[English source not present in the checked-in Bencraft/Joly files.]}
}
{
}

\Verse{8}
{
    \zR{匪曰无 故呻吟，亦长歌当哭之意耳。}
}
{
    \eR{[English source not present in the checked-in Bencraft/Joly files.]}
}
{
}

\Verse{9}
{
    \zR{悲时序之递嬗兮，又属清秋。}
}
{
    \eR{[English source not present in the checked-in Bencraft/Joly files.]}
}
{
}

\Verse{10}
{
    \zR{感遭家之不造兮，独处离愁。}
}
{
    \eR{[English source not present in the checked-in Bencraft/Joly files.]}
}
{
}

\Verse{11}
{
    \zR{北堂有萱兮，何以 忘忧？}
}
{
    \eR{[English source not present in the checked-in Bencraft/Joly files.]}
}
{
}

\Verse{12}
{
    \zR{无以解忧兮，我心咻咻。}
}
{
    \eR{[English source not present in the checked-in Bencraft/Joly files.]}
}
{
}

\Verse{13}
{
    \zR{云凭凭兮秋风酸，步中庭兮霜叶干。}
}
{
    \eR{[English source not present in the checked-in Bencraft/Joly files.]}
}
{
}

\Verse{14}
{
    \zR{何去何从兮失我故欢，静言思之兮恻肺肝。}
}
{
    \eR{[English source not present in the checked-in Bencraft/Joly files.]}
}
{
}

\Verse{15}
{
    \zR{惟鲔有潭兮，惟鹤有梁。}
}
{
    \eR{[English source not present in the checked-in Bencraft/Joly files.]}
}
{
}

\Verse{16}
{
    \zR{鳞甲潜伏兮，羽毛何长!搔首问兮茫茫，高天厚地兮， 谁知余之永伤？}
}
{
    \eR{[English source not present in the checked-in Bencraft/Joly files.]}
}
{
}

\Verse{17}
{
    \zR{银河耿耿兮寒气侵，月色横斜兮玉漏沉。}
}
{
    \eR{[English source not present in the checked-in Bencraft/Joly files.]}
}
{
}

\Verse{18}
{
    \zR{忧心炳炳兮发我哀吟。}
}
{
    \eR{[English source not present in the checked-in Bencraft/Joly files.]}
}
{
}

\Verse{19}
{
    \zR{吟复吟兮寄我 知音。}
}
{
    \eR{[English source not present in the checked-in Bencraft/Joly files.]}
}
{
}

\Verse{20}
{
    \zR{黛玉看了，不胜伤感。}
}
{
    \eR{[English source not present in the checked-in Bencraft/Joly files.]}
}
{
}

\Verse{21}
{
    \zR{又想：“宝姐姐不寄与别人，单寄与我，也是‘惺惺惜惺惺’ 的意思。”}
}
{
    \eR{[English source not present in the checked-in Bencraft/Joly files.]}
}
{
}

\Verse{22}
{
    \zR{正在沉吟，只听见外面有人说道：“林姐姐在家里呢么？”}
}
{
    \eR{[English source not present in the checked-in Bencraft/Joly files.]}
}
{
}

\Verse{23}
{
    \zR{黛玉一面把 宝钗的书叠起，口内便答应道：“是谁？”}
}
{
    \eR{[English source not present in the checked-in Bencraft/Joly files.]}
}
{
}

\Verse{24}
{
    \zR{正问着，早见几个人进来，却是探春、 湘云、李纹、李绮。}
}
{
    \eR{[English source not present in the checked-in Bencraft/Joly files.]}
}
{
}

\Verse{25}
{
    \zR{彼此问了好，雪雁倒上茶来，大家喝了，说些闲话。}
}
{
    \eR{[English source not present in the checked-in Bencraft/Joly files.]}
}
{
}

\Verse{26}
{
    \zR{因想起前 年的“菊花诗”来，黛玉便道：“宝姐姐自从挪出去，来了两遭，如今索性有事也 不来了，真真奇怪。}
}
{
    \eR{[English source not present in the checked-in Bencraft/Joly files.]}
}
{
}

\Verse{27}
{
    \zR{我看他终久还来我们这里不来！”}
}
{
    \eR{[English source not present in the checked-in Bencraft/Joly files.]}
}
{
}

\Verse{28}
{
    \zR{探春微笑道：“怎么不来？}
}
{
    \eR{[English source not present in the checked-in Bencraft/Joly files.]}
}
{
}

\Verse{29}
{
    \zR{横竖要来的。}
}
{
    \eR{[English source not present in the checked-in Bencraft/Joly files.]}
}
{
}

\Verse{30}
{
    \zR{如今是他们尊嫂有些脾气，姨妈上了年纪的人，又兼有薛大哥的事， 自然得宝姐姐照料一切。}
}
{
    \eR{[English source not present in the checked-in Bencraft/Joly files.]}
}
{
}

\Verse{31}
{
    \zR{那里还比得先前有工夫呢？”}
}
{
    \eR{[English source not present in the checked-in Bencraft/Joly files.]}
}
{
}

\Verse{32}
{
    \zR{正说着，忽听得唿喇喇一片风声，吹了好些落叶打在窗纸上。}
}
{
    \eR{[English source not present in the checked-in Bencraft/Joly files.]}
}
{
}

\Verse{33}
{
    \zR{停了一回儿，又 透过一阵清香来。}
}
{
    \eR{[English source not present in the checked-in Bencraft/Joly files.]}
}
{
}

\Verse{34}
{
    \zR{众人闻着，都说道：“这是何处来的香风？}
}
{
    \eR{[English source not present in the checked-in Bencraft/Joly files.]}
}
{
}

\Verse{35}
{
    \zR{这像什么香？”}
}
{
    \eR{[English source not present in the checked-in Bencraft/Joly files.]}
}
{
}

\Verse{36}
{
    \zR{黛玉 道：“好像木樨香。”}
}
{
    \eR{[English source not present in the checked-in Bencraft/Joly files.]}
}
{
}

\Verse{37}
{
    \zR{探春笑道：“林姐姐终不脱南边人的话。}
}
{
    \eR{[English source not present in the checked-in Bencraft/Joly files.]}
}
{
}

\Verse{38}
{
    \zR{这大九月里的，那 里还有桂花呢？”}
}
{
    \eR{[English source not present in the checked-in Bencraft/Joly files.]}
}
{
}

\Verse{39}
{
    \zR{黛玉笑道：“原是啊!不然，怎么不竟说‘是’桂花香，只说似 乎‘像’呢？”}
}
{
    \eR{[English source not present in the checked-in Bencraft/Joly files.]}
}
{
}

\Verse{40}
{
    \zR{湘云道：“三姐姐，你也别说。}
}
{
    \eR{[English source not present in the checked-in Bencraft/Joly files.]}
}
{
}

\Verse{41}
{
    \zR{你可记得‘十里荷花，三秋桂子’？}
}
{
    \eR{[English source not present in the checked-in Bencraft/Joly files.]}
}
{
}

\Verse{42}
{
    \zR{在南边正是晚桂开的时候了，你只没有见过罢了。}
}
{
    \eR{[English source not present in the checked-in Bencraft/Joly files.]}
}
{
}

\Verse{43}
{
    \zR{等你明日到南边去的时候，你自 然也就知道了。”}
}
{
    \eR{[English source not present in the checked-in Bencraft/Joly files.]}
}
{
}

\Verse{44}
{
    \zR{探春笑道：“我有什么事到南边去？}
}
{
    \eR{[English source not present in the checked-in Bencraft/Joly files.]}
}
{
}

\Verse{45}
{
    \zR{况且这个也是我早知道的， 不用你们说嘴。”}
}
{
    \eR{[English source not present in the checked-in Bencraft/Joly files.]}
}
{
}

\Verse{46}
{
    \zR{李纹、李绮只抿着嘴儿笑。}
}
{
    \eR{[English source not present in the checked-in Bencraft/Joly files.]}
}
{
}

\Verse{47}
{
    \zR{黛玉道：“妹妹，这可说不齐。}
}
{
    \eR{[English source not present in the checked-in Bencraft/Joly files.]}
}
{
}

\Verse{48}
{
    \zR{俗语 说：‘人是地行仙。’}
}
{
    \eR{[English source not present in the checked-in Bencraft/Joly files.]}
}
{
}

\Verse{49}
{
    \zR{今日在这里，明日就不知在那里。}
}
{
    \eR{[English source not present in the checked-in Bencraft/Joly files.]}
}
{
}

\Verse{50}
{
    \zR{譬如我原是南边人，怎么 到了这里呢？”}
}
{
    \eR{[English source not present in the checked-in Bencraft/Joly files.]}
}
{
}

\Verse{51}
{
    \zR{湘云拍着手笑道：“今儿三姐姐可叫林姐姐问住了。}
}
{
    \eR{[English source not present in the checked-in Bencraft/Joly files.]}
}
{
}

\Verse{52}
{
    \zR{不但林姐姐是 南边人到这里，就是我们这几个人就不同：也有本来是北边的；也有根子是南边， 生长在北边的；也有生长在南边，到这北边的。}
}
{
    \eR{[English source not present in the checked-in Bencraft/Joly files.]}
}
{
}

\Verse{53}
{
    \zR{今儿大家都凑在一处，可见人总有 一个定数。}
}
{
    \eR{[English source not present in the checked-in Bencraft/Joly files.]}
}
{
}

\Verse{54}
{
    \zR{大凡地和人，总是各自有缘分的。”}
}
{
    \eR{[English source not present in the checked-in Bencraft/Joly files.]}
}
{
}

\Verse{55}
{
    \zR{众人听了都点头，探春也只是笑。}
}
{
    \eR{[English source not present in the checked-in Bencraft/Joly files.]}
}
{
}

\Verse{56}
{
    \zR{又说了一会子闲话儿，大家散出。}
}
{
    \eR{[English source not present in the checked-in Bencraft/Joly files.]}
}
{
}

\Verse{57}
{
    \zR{黛玉送至门口，大家都说：“你身上才好些，别 出来了，看着了风。”}
}
{
    \eR{[English source not present in the checked-in Bencraft/Joly files.]}
}
{
}

\Verse{58}
{
    \zR{于是黛玉一面说着话儿，一面站在门口，又与四人殷勤了几句，便看着他们出 院去了。}
}
{
    \eR{[English source not present in the checked-in Bencraft/Joly files.]}
}
{
}

\Verse{59}
{
    \zR{进来坐着，看看已是林鸟归山，夕阳西坠。}
}
{
    \eR{[English source not present in the checked-in Bencraft/Joly files.]}
}
{
}

\Verse{60}
{
    \zR{因史湘云说起南边的话，便想 着：“父母若在，南边的景致，春花秋月，水秀山明，二十四桥，六朝遗迹。}
}
{
    \eR{[English source not present in the checked-in Bencraft/Joly files.]}
}
{
}

\Verse{61}
{
    \zR{不少 下人伏侍，诸事可以任意，言语亦可不避。}
}
{
    \eR{[English source not present in the checked-in Bencraft/Joly files.]}
}
{
}

\Verse{62}
{
    \zR{香车画舫，红杏青帘，惟我独尊。}
}
{
    \eR{[English source not present in the checked-in Bencraft/Joly files.]}
}
{
}

\Verse{63}
{
    \zR{今日 寄人篱下，纵有许多照应，自己无处不要留心。}
}
{
    \eR{[English source not present in the checked-in Bencraft/Joly files.]}
}
{
}

\Verse{64}
{
    \zR{不知前生作了什么罪孽，今生这样 孤凄!真是李后主说的，‘此间日中只以眼泪洗面’矣！”}
}
{
    \eR{[English source not present in the checked-in Bencraft/Joly files.]}
}
{
}

\Verse{65}
{
    \zR{一面思想，不知不觉神 往那里去了。}
}
{
    \eR{[English source not present in the checked-in Bencraft/Joly files.]}
}
{
}

\Verse{66}
{
    \zR{紫鹃走来，看见这样光景，想着必是因刚才说起南边北边的话来，一时触着黛 玉的心事了。}
}
{
    \eR{[English source not present in the checked-in Bencraft/Joly files.]}
}
{
}

\Verse{67}
{
    \zR{便问道：“姑娘们来说了半天话，想来姑娘又劳了神了。}
}
{
    \eR{[English source not present in the checked-in Bencraft/Joly files.]}
}
{
}

\Verse{68}
{
    \zR{刚才我叫雪 雁告诉厨房里，给姑娘作了一碗火肉白菜汤，加了一点儿虾米儿，配了点青笋紫菜， 姑娘想着好么？”}
}
{
    \eR{[English source not present in the checked-in Bencraft/Joly files.]}
}
{
}

\Verse{69}
{
    \zR{黛玉道：“也罢了。”}
}
{
    \eR{[English source not present in the checked-in Bencraft/Joly files.]}
}
{
}

\Verse{70}
{
    \zR{紫鹃道：“还熬了一点江米粥。”}
}
{
    \eR{[English source not present in the checked-in Bencraft/Joly files.]}
}
{
}

\Verse{71}
{
    \zR{黛玉点 点头儿，又说道：“那粥得你们两个自己熬了，不用他们厨房里熬才是。”}
}
{
    \eR{[English source not present in the checked-in Bencraft/Joly files.]}
}
{
}

\Verse{72}
{
    \zR{紫鹃道： “我也怕厨房里弄的不干净，我们自己熬呢。}
}
{
    \eR{[English source not present in the checked-in Bencraft/Joly files.]}
}
{
}

\Verse{73}
{
    \zR{就是那汤，我也告诉雪雁合柳嫂儿说 了，要弄干净着。}
}
{
    \eR{[English source not present in the checked-in Bencraft/Joly files.]}
}
{
}

\Verse{74}
{
    \zR{柳嫂儿说了：他打点妥当，拿到他屋里，叫他们五儿瞅着炖呢。”}
}
{
    \eR{[English source not present in the checked-in Bencraft/Joly files.]}
}
{
}

\Verse{75}
{
    \zR{黛玉道：“我倒不是嫌人家腌？。}
}
{
    \eR{[English source not present in the checked-in Bencraft/Joly files.]}
}
{
}

\Verse{76}
{
    \zR{只是病了好些日子，不周不备，都是人家，这会 子又汤儿粥儿的调度，未免惹人厌烦。”}
}
{
    \eR{[English source not present in the checked-in Bencraft/Joly files.]}
}
{
}

\Verse{77}
{
    \zR{说着，眼圈儿又红了。}
}
{
    \eR{[English source not present in the checked-in Bencraft/Joly files.]}
}
{
}

\Verse{78}
{
    \zR{紫鹃道：“姑娘这 话也是多想。}
}
{
    \eR{[English source not present in the checked-in Bencraft/Joly files.]}
}
{
}

\Verse{79}
{
    \zR{姑娘是老太太的外孙女儿，又是老太太心坎儿上的。}
}
{
    \eR{[English source not present in the checked-in Bencraft/Joly files.]}
}
{
}

\Verse{80}
{
    \zR{别人求其在姑娘 跟前讨好儿还不能呢，那里有抱怨的？”}
}
{
    \eR{[English source not present in the checked-in Bencraft/Joly files.]}
}
{
}

\Verse{81}
{
    \zR{黛玉点点头儿。}
}
{
    \eR{[English source not present in the checked-in Bencraft/Joly files.]}
}
{
}

\Verse{82}
{
    \zR{因又问道：“你才说的五 儿，不是那日合宝二爷那边的芳官在一处的那个女孩儿？”}
}
{
    \eR{[English source not present in the checked-in Bencraft/Joly files.]}
}
{
}

\Verse{83}
{
    \zR{紫鹃道：“就是他。”}
}
{
    \eR{[English source not present in the checked-in Bencraft/Joly files.]}
}
{
}

\Verse{84}
{
    \zR{黛玉道：“不听见说要进来么？”}
}
{
    \eR{[English source not present in the checked-in Bencraft/Joly files.]}
}
{
}

\Verse{85}
{
    \zR{紫鹃道：“可不是，因为病了一场。}
}
{
    \eR{[English source not present in the checked-in Bencraft/Joly files.]}
}
{
}

\Verse{86}
{
    \zR{后来好了， 才要进来，正是晴雯他们闹出事来的时候，也就耽搁住了。”}
}
{
    \eR{[English source not present in the checked-in Bencraft/Joly files.]}
}
{
}

\Verse{87}
{
    \zR{黛玉道：“我看那丫 头倒也还头脸儿干净。”}
}
{
    \eR{[English source not present in the checked-in Bencraft/Joly files.]}
}
{
}

\Verse{88}
{
    \zR{说着，外头婆子送了汤来。}
}
{
    \eR{[English source not present in the checked-in Bencraft/Joly files.]}
}
{
}

\Verse{89}
{
    \zR{雪雁出来接时，那婆子说道： “柳嫂儿叫回姑娘：这是他们五儿作的，没敢在大厨房里作，怕姑娘嫌腌？。”}
}
{
    \eR{[English source not present in the checked-in Bencraft/Joly files.]}
}
{
}

\Verse{90}
{
    \zR{雪 雁答应着，接了进来。}
}
{
    \eR{[English source not present in the checked-in Bencraft/Joly files.]}
}
{
}

\Verse{91}
{
    \zR{黛玉在屋里，已听见了，吩咐雪雁：“告诉那老婆子回去说， 叫他费心。”}
}
{
    \eR{[English source not present in the checked-in Bencraft/Joly files.]}
}
{
}

\Verse{92}
{
    \zR{雪雁出来说了，老婆子自去。}
}
{
    \eR{[English source not present in the checked-in Bencraft/Joly files.]}
}
{
}

\Verse{93}
{
    \zR{这里雪雁将黛玉的碗箸安放在小几儿上， 因问黛玉道：“还有咱们南来的五香大头菜，拌些麻油、醋，可好么？”}
}
{
    \eR{[English source not present in the checked-in Bencraft/Joly files.]}
}
{
}

\Verse{94}
{
    \zR{黛玉道： “也使得，只不必累坠了。”}
}
{
    \eR{[English source not present in the checked-in Bencraft/Joly files.]}
}
{
}

\Verse{95}
{
    \zR{一面盛上粥来。}
}
{
    \eR{[English source not present in the checked-in Bencraft/Joly files.]}
}
{
}

\Verse{96}
{
    \zR{黛玉吃了半碗，用羹匙舀了两口汤喝， 就搁下了。}
}
{
    \eR{[English source not present in the checked-in Bencraft/Joly files.]}
}
{
}

\Verse{97}
{
    \zR{两个丫鬟撤下来了，拭净了小几，端下去，又换上一张常放的小几。}
}
{
    \eR{[English source not present in the checked-in Bencraft/Joly files.]}
}
{
}

\Verse{98}
{
    \zR{黛 玉漱了口，盥了手，便道：“紫鹃，添了香了没有？”}
}
{
    \eR{[English source not present in the checked-in Bencraft/Joly files.]}
}
{
}

\Verse{99}
{
    \zR{紫鹃道：“就添去。”}
}
{
    \eR{[English source not present in the checked-in Bencraft/Joly files.]}
}
{
}

\Verse{100}
{
    \zR{黛玉 道：“你们就把那汤合粥吃了罢，味儿还好，且是干净。}
}
{
    \eR{[English source not present in the checked-in Bencraft/Joly files.]}
}
{
}

\Verse{101}
{
    \zR{待我自己添香罢。”}
}
{
    \eR{[English source not present in the checked-in Bencraft/Joly files.]}
}
{
}

\Verse{102}
{
    \zR{两个 人答应了，在外间自吃去了。}
}
{
    \eR{[English source not present in the checked-in Bencraft/Joly files.]}
}
{
}

\Verse{103}
{
    \zR{这里黛玉添了香，自己坐着，才要拿本书看，只听得园内的风自西边直透到东 边，穿过树枝，都在那里唏？}
}
{
    \eR{[English source not present in the checked-in Bencraft/Joly files.]}
}
{
}

\Verse{104}
{
    \zR{哗喇不住的响。}
}
{
    \eR{[English source not present in the checked-in Bencraft/Joly files.]}
}
{
}

\Verse{105}
{
    \zR{一会儿，檐下的铁马也只管叮叮当当 的乱敲起来。}
}
{
    \eR{[English source not present in the checked-in Bencraft/Joly files.]}
}
{
}

\Verse{106}
{
    \zR{一时雪雁先吃完了，进来伺候。}
}
{
    \eR{[English source not present in the checked-in Bencraft/Joly files.]}
}
{
}

\Verse{107}
{
    \zR{黛玉便问道：“天气冷了，我前日叫 你们把那些小毛儿衣裳晾晾，可曾晾过没有？”}
}
{
    \eR{[English source not present in the checked-in Bencraft/Joly files.]}
}
{
}

\Verse{108}
{
    \zR{雪雁道：“都晾过了。”}
}
{
    \eR{[English source not present in the checked-in Bencraft/Joly files.]}
}
{
}

\Verse{109}
{
    \zR{黛玉道： “你拿一件来我披披。”}
}
{
    \eR{[English source not present in the checked-in Bencraft/Joly files.]}
}
{
}

\Verse{110}
{
    \zR{雪雁走去，将一包小毛衣裳抱来，打开毡包，给黛玉自拣。}
}
{
    \eR{[English source not present in the checked-in Bencraft/Joly files.]}
}
{
}

\Verse{111}
{
    \zR{只见内中夹着个绢包儿。}
}
{
    \eR{[English source not present in the checked-in Bencraft/Joly files.]}
}
{
}

\Verse{112}
{
    \zR{黛玉伸手拿起，打开看时，却是宝玉病时送来的旧绢子， 自己题的诗，上面泪痕犹在。}
}
{
    \eR{[English source not present in the checked-in Bencraft/Joly files.]}
}
{
}

\Verse{113}
{
    \zR{里头却包着那剪破了的香囊、扇袋并宝玉通灵玉上的 穗子。}
}
{
    \eR{[English source not present in the checked-in Bencraft/Joly files.]}
}
{
}

\Verse{114}
{
    \zR{原来晾衣裳时从箱中检出，紫鹃恐怕遗失了，遂夹在这毡包里的。}
}
{
    \eR{[English source not present in the checked-in Bencraft/Joly files.]}
}
{
}

\Verse{115}
{
    \zR{这黛玉不 看则已，看了时，也不说穿那一件衣裳，手里只拿着那两方手帕，呆呆的看那旧诗。}
}
{
    \eR{[English source not present in the checked-in Bencraft/Joly files.]}
}
{
}

\Verse{116}
{
    \zR{看了一回，不觉得簌簌泪下。}
}
{
    \eR{[English source not present in the checked-in Bencraft/Joly files.]}
}
{
}

\Verse{117}
{
    \zR{紫鹃刚从外间进来，只见雪雁正捧着一毡包衣裳，在傍边呆立，小几上却搁着 剪破了的香囊和两三截儿扇袋并那铰拆了的穗子。}
}
{
    \eR{[English source not present in the checked-in Bencraft/Joly files.]}
}
{
}

\Verse{118}
{
    \zR{黛玉手中却拿着两方旧帕子，上 边写着字迹，在那里对着滴泪呢。}
}
{
    \eR{[English source not present in the checked-in Bencraft/Joly files.]}
}
{
}

\Verse{119}
{
    \zR{正是： 失意人逢失意事，新啼痕间旧啼痕。}
}
{
    \eR{[English source not present in the checked-in Bencraft/Joly files.]}
}
{
}

\Verse{120}
{
    \zR{紫鹃见了这样，知是他触物伤情，感怀旧事，料道劝也无益，只得笑着道：“姑娘， 还看那些东西作什么？}
}
{
    \eR{[English source not present in the checked-in Bencraft/Joly files.]}
}
{
}

\Verse{121}
{
    \zR{那都是那几年宝二爷和姑娘小时，一时好了，一时恼了，闹 出来的笑话儿。}
}
{
    \eR{[English source not present in the checked-in Bencraft/Joly files.]}
}
{
}

\Verse{122}
{
    \zR{要像如今这样厮抬厮敬的，那里能把这些东西白遭塌了呢。”}
}
{
    \eR{[English source not present in the checked-in Bencraft/Joly files.]}
}
{
}

\Verse{123}
{
    \zR{紫鹃 这话原给黛玉开心，不料这几句话更提起黛玉初来时和宝玉的旧事来，一发珠泪连 绵起来。}
}
{
    \eR{[English source not present in the checked-in Bencraft/Joly files.]}
}
{
}

\Verse{124}
{
    \zR{紫鹃又劝道：“雪雁这里等着呢，姑娘披上一件罢。”}
}
{
    \eR{[English source not present in the checked-in Bencraft/Joly files.]}
}
{
}

\Verse{125}
{
    \zR{那黛玉才把手帕撂 下。}
}
{
    \eR{[English source not present in the checked-in Bencraft/Joly files.]}
}
{
}

\Verse{126}
{
    \zR{紫鹃连忙拾起，将香袋等物包起拿开。}
}
{
    \eR{[English source not present in the checked-in Bencraft/Joly files.]}
}
{
}

\Verse{127}
{
    \zR{这黛玉方披了一件皮衣，自己闷闷的走 到外间来坐下。}
}
{
    \eR{[English source not present in the checked-in Bencraft/Joly files.]}
}
{
}

\Verse{128}
{
    \zR{回头看见案上宝钗的诗启尚未收好，又拿出来瞧了两遍，叹道：“境 遇不同，伤心则一。}
}
{
    \eR{[English source not present in the checked-in Bencraft/Joly files.]}
}
{
}

\Verse{129}
{
    \zR{不免也赋四章，翻入琴谱，可弹可歌，明日写出来寄去，以当 和作。”}
}
{
    \eR{[English source not present in the checked-in Bencraft/Joly files.]}
}
{
}

\Verse{130}
{
    \zR{便叫雪雁将外边桌上笔砚拿来，濡墨挥毫，赋成四叠。}
}
{
    \eR{[English source not present in the checked-in Bencraft/Joly files.]}
}
{
}

\Verse{131}
{
    \zR{又将琴谱翻出，借 他《猗兰》《思贤》两操，合成音韵，与自己做的配齐了，然后写出，以备送与宝 钗。}
}
{
    \eR{[English source not present in the checked-in Bencraft/Joly files.]}
}
{
}

\Verse{132}
{
    \zR{又即叫雪雁向箱中将自己带来的短琴拿出，调上弦，又操演了指法。}
}
{
    \eR{[English source not present in the checked-in Bencraft/Joly files.]}
}
{
}

\Verse{133}
{
    \zR{黛玉本是 个绝顶聪明人，又在南边学过几时，虽是手生，到底一理就熟。}
}
{
    \eR{[English source not present in the checked-in Bencraft/Joly files.]}
}
{
}

\Verse{134}
{
    \zR{抚了一番，夜已深 了，便叫紫鹃收拾睡觉，不提。}
}
{
    \eR{[English source not present in the checked-in Bencraft/Joly files.]}
}
{
}

\Verse{135}
{
    \zR{却说宝玉这日起来，梳洗了，带着焙茗正往书房中来，只见墨雨笑嘻嘻的跑来， 迎头说道：“二爷今日便宜了。}
}
{
    \eR{[English source not present in the checked-in Bencraft/Joly files.]}
}
{
}

\Verse{136}
{
    \zR{太爷不在书房里，都放了学了。”}
}
{
    \eR{[English source not present in the checked-in Bencraft/Joly files.]}
}
{
}

\Verse{137}
{
    \zR{宝玉道：“当真 的么？”}
}
{
    \eR{[English source not present in the checked-in Bencraft/Joly files.]}
}
{
}

\Verse{138}
{
    \zR{墨雨道：“二爷不信，那不是三爷和兰哥来了？”}
}
{
    \eR{[English source not present in the checked-in Bencraft/Joly files.]}
}
{
}

\Verse{139}
{
    \zR{宝玉看时，只见贾环贾 兰跟着小厮们，两个笑嘻嘻的，嘴里咭咭呱呱不知说些什么，迎头来了。}
}
{
    \eR{[English source not present in the checked-in Bencraft/Joly files.]}
}
{
}

\Verse{140}
{
    \zR{见了宝玉， 都垂手站住。}
}
{
    \eR{[English source not present in the checked-in Bencraft/Joly files.]}
}
{
}

\Verse{141}
{
    \zR{宝玉问道：“你们两个怎么就回来了？”}
}
{
    \eR{[English source not present in the checked-in Bencraft/Joly files.]}
}
{
}

\Verse{142}
{
    \zR{贾环道：“今日太爷有事， 说是放一天学，明儿再去呢。”}
}
{
    \eR{[English source not present in the checked-in Bencraft/Joly files.]}
}
{
}

\Verse{143}
{
    \zR{宝玉听了，方回身到贾母贾政处去禀明了，然后回 到怡红院中。}
}
{
    \eR{[English source not present in the checked-in Bencraft/Joly files.]}
}
{
}

\Verse{144}
{
    \zR{袭人问道：“怎么又回来了？”}
}
{
    \eR{[English source not present in the checked-in Bencraft/Joly files.]}
}
{
}

\Verse{145}
{
    \zR{宝玉告诉了他。}
}
{
    \eR{[English source not present in the checked-in Bencraft/Joly files.]}
}
{
}

\Verse{146}
{
    \zR{只坐了一坐儿，便往 外走，袭人道：“往那里去，这样忙法？}
}
{
    \eR{[English source not present in the checked-in Bencraft/Joly files.]}
}
{
}

\Verse{147}
{
    \zR{就放了学，依我说，也该养养神儿了。”}
}
{
    \eR{[English source not present in the checked-in Bencraft/Joly files.]}
}
{
}

\Verse{148}
{
    \zR{宝玉站住脚，低了头，说道：“你的话也是，但是好容易放一天学，还不散散去。}
}
{
    \eR{[English source not present in the checked-in Bencraft/Joly files.]}
}
{
}

\Verse{149}
{
    \zR{你也该可怜我些儿了。”}
}
{
    \eR{[English source not present in the checked-in Bencraft/Joly files.]}
}
{
}

\Verse{150}
{
    \zR{袭人见说的可怜，笑道：“由爷去罢。”}
}
{
    \eR{[English source not present in the checked-in Bencraft/Joly files.]}
}
{
}

\Verse{151}
{
    \zR{正说着，端了饭 来，宝玉也没法儿，只得且吃饭。}
}
{
    \eR{[English source not present in the checked-in Bencraft/Joly files.]}
}
{
}

\Verse{152}
{
    \zR{三口两口忙忙的吃完，漱了口，一溜烟往黛玉房 中去了。}
}
{
    \eR{[English source not present in the checked-in Bencraft/Joly files.]}
}
{
}

\Verse{153}
{
    \zR{走到门口，只见雪雁在院中晾绢子呢。}
}
{
    \eR{[English source not present in the checked-in Bencraft/Joly files.]}
}
{
}

\Verse{154}
{
    \zR{宝玉因问：“姑娘吃了饭了么？”}
}
{
    \eR{[English source not present in the checked-in Bencraft/Joly files.]}
}
{
}

\Verse{155}
{
    \zR{雪雁 道：“早起喝了半碗粥，懒怠吃饭，这时候打盹儿呢。}
}
{
    \eR{[English source not present in the checked-in Bencraft/Joly files.]}
}
{
}

\Verse{156}
{
    \zR{二爷且到别处走走，回来再 来罢。”}
}
{
    \eR{[English source not present in the checked-in Bencraft/Joly files.]}
}
{
}

\Verse{157}
{
    \zR{宝玉只得回来。}
}
{
    \eR{[English source not present in the checked-in Bencraft/Joly files.]}
}
{
}

\Verse{158}
{
    \zR{无处可去，忽然想起惜春有好几天没见，便信步走到蓼风 轩来。}
}
{
    \eR{[English source not present in the checked-in Bencraft/Joly files.]}
}
{
}

\Verse{159}
{
    \zR{刚到窗下，只见静悄悄一无人声，宝玉打量他也睡午觉，不便进去。}
}
{
    \eR{[English source not present in the checked-in Bencraft/Joly files.]}
}
{
}

\Verse{160}
{
    \zR{才要走 时，只听屋里微微一响，不知何声；宝玉站住再听，半日，又“拍”的一响。}
}
{
    \eR{[English source not present in the checked-in Bencraft/Joly files.]}
}
{
}

\Verse{161}
{
    \zR{宝玉 还未听出，只见一个人道：“你在这里下了一个子儿，那里你不应么？”}
}
{
    \eR{[English source not present in the checked-in Bencraft/Joly files.]}
}
{
}

\Verse{162}
{
    \zR{宝玉方知 是下棋呢。}
}
{
    \eR{[English source not present in the checked-in Bencraft/Joly files.]}
}
{
}

\Verse{163}
{
    \zR{但只急切听不出这个人的语音是谁。}
}
{
    \eR{[English source not present in the checked-in Bencraft/Joly files.]}
}
{
}

\Verse{164}
{
    \zR{底下方听见惜春道：“怕什么？}
}
{
    \eR{[English source not present in the checked-in Bencraft/Joly files.]}
}
{
}

\Verse{165}
{
    \zR{你 这么一吃我，我这么一应；你又这么吃，我又这么应：还缓着一着儿呢，终久连的 上。”}
}
{
    \eR{[English source not present in the checked-in Bencraft/Joly files.]}
}
{
}

\Verse{166}
{
    \zR{那一个又道：“我要这么一吃呢？”}
}
{
    \eR{[English source not present in the checked-in Bencraft/Joly files.]}
}
{
}

\Verse{167}
{
    \zR{惜春道：“啊嗄，还有一着反扑在里头 呢，我倒没防备。”}
}
{
    \eR{[English source not present in the checked-in Bencraft/Joly files.]}
}
{
}

\Verse{168}
{
    \zR{宝玉听了听那一个声音很熟，却不是他们姊妹，料着惜春屋里 也没外人，轻轻的掀帘进去。}
}
{
    \eR{[English source not present in the checked-in Bencraft/Joly files.]}
}
{
}

\Verse{169}
{
    \zR{看时不是别人，却是那栊翠庵的槛外人妙玉。}
}
{
    \eR{[English source not present in the checked-in Bencraft/Joly files.]}
}
{
}

\Verse{170}
{
    \zR{这宝玉 见是妙玉，不敢惊动。}
}
{
    \eR{[English source not present in the checked-in Bencraft/Joly files.]}
}
{
}

\Verse{171}
{
    \zR{妙玉和惜春正在凝思之际，也没理会。}
}
{
    \eR{[English source not present in the checked-in Bencraft/Joly files.]}
}
{
}

\Verse{172}
{
    \zR{宝玉却站在旁边，看 他两个的手段。}
}
{
    \eR{[English source not present in the checked-in Bencraft/Joly files.]}
}
{
}

\Verse{173}
{
    \zR{只见妙玉低着头，问惜春道：“你这个畸角儿不要了么？”}
}
{
    \eR{[English source not present in the checked-in Bencraft/Joly files.]}
}
{
}

\Verse{174}
{
    \zR{惜春道： “怎么不要？}
}
{
    \eR{[English source not present in the checked-in Bencraft/Joly files.]}
}
{
}

\Verse{175}
{
    \zR{你那里头都是死子儿，我怕什么？”}
}
{
    \eR{[English source not present in the checked-in Bencraft/Joly files.]}
}
{
}

\Verse{176}
{
    \zR{妙玉道：“且别说满话，试试看。”}
}
{
    \eR{[English source not present in the checked-in Bencraft/Joly files.]}
}
{
}

\Verse{177}
{
    \zR{惜春道：“我便打了起来，看你怎么着。”}
}
{
    \eR{[English source not present in the checked-in Bencraft/Joly files.]}
}
{
}

\Verse{178}
{
    \zR{妙玉却微微笑着，把边上子一接，却搭 转一吃，把惜春的一个角儿都打起来了，笑着说道：“这叫做‘倒脱靴势’。”}
}
{
    \eR{[English source not present in the checked-in Bencraft/Joly files.]}
}
{
}

\Verse{179}
{
    \zR{惜春尚未答言，宝玉在旁情不自禁，哈哈一笑，把两个人都唬了一大跳。}
}
{
    \eR{[English source not present in the checked-in Bencraft/Joly files.]}
}
{
}

\Verse{180}
{
    \zR{惜春 道：“你这是怎么说？}
}
{
    \eR{[English source not present in the checked-in Bencraft/Joly files.]}
}
{
}

\Verse{181}
{
    \zR{进来也不言语。}
}
{
    \eR{[English source not present in the checked-in Bencraft/Joly files.]}
}
{
}

\Verse{182}
{
    \zR{这么使促狭唬人!你多早晚进来的？”}
}
{
    \eR{[English source not present in the checked-in Bencraft/Joly files.]}
}
{
}

\Verse{183}
{
    \zR{宝玉道： “我头里就进来了，看着你们两个争这个畸角儿。”}
}
{
    \eR{[English source not present in the checked-in Bencraft/Joly files.]}
}
{
}

\Verse{184}
{
    \zR{说着，一面与妙玉施礼，一面 又笑问道：“妙公轻易不出禅关，今日何缘下凡一走？”}
}
{
    \eR{[English source not present in the checked-in Bencraft/Joly files.]}
}
{
}

\Verse{185}
{
    \zR{妙玉听了，忽然把脸一红， 也不答言，低了头自看那棋。}
}
{
    \eR{[English source not present in the checked-in Bencraft/Joly files.]}
}
{
}

\Verse{186}
{
    \zR{宝玉自觉造次，连忙陪笑道：“倒是出家人比不得我 们在家的俗人。}
}
{
    \eR{[English source not present in the checked-in Bencraft/Joly files.]}
}
{
}

\Verse{187}
{
    \zR{头一件，心是静的。}
}
{
    \eR{[English source not present in the checked-in Bencraft/Joly files.]}
}
{
}

\Verse{188}
{
    \zR{静则灵，灵则慧。”}
}
{
    \eR{[English source not present in the checked-in Bencraft/Joly files.]}
}
{
}

\Verse{189}
{
    \zR{宝玉尚未说完，只见妙玉 微微的把眼一抬，看了宝玉一眼，复又低下头去，那脸上的颜色渐渐的红晕起来。}
}
{
    \eR{[English source not present in the checked-in Bencraft/Joly files.]}
}
{
}

\Verse{190}
{
    \zR{宝玉见他不理，只得讪讪的旁边坐了。}
}
{
    \eR{[English source not present in the checked-in Bencraft/Joly files.]}
}
{
}

\Verse{191}
{
    \zR{惜春还要下子，妙玉半日说道：“再下罢。”}
}
{
    \eR{[English source not present in the checked-in Bencraft/Joly files.]}
}
{
}

\Verse{192}
{
    \zR{便起身理理衣裳，重新坐下，痴 痴的问着宝玉道：“你从何处来？”}
}
{
    \eR{[English source not present in the checked-in Bencraft/Joly files.]}
}
{
}

\Verse{193}
{
    \zR{宝玉巴不得这一声，好解释前头的话，忽又想 道：“或是妙玉的机锋？”}
}
{
    \eR{[English source not present in the checked-in Bencraft/Joly files.]}
}
{
}

\Verse{194}
{
    \zR{转红了脸，答应不出来。}
}
{
    \eR{[English source not present in the checked-in Bencraft/Joly files.]}
}
{
}

\Verse{195}
{
    \zR{妙玉微微一笑，自合惜春说话。}
}
{
    \eR{[English source not present in the checked-in Bencraft/Joly files.]}
}
{
}

\Verse{196}
{
    \zR{惜春也笑道：“二哥哥，这什么难答的？}
}
{
    \eR{[English source not present in the checked-in Bencraft/Joly files.]}
}
{
}

\Verse{197}
{
    \zR{你没有听见人家常说的，‘从来处来’么？}
}
{
    \eR{[English source not present in the checked-in Bencraft/Joly files.]}
}
{
}

\Verse{198}
{
    \zR{这也值得把脸红了，见了生人的似的。”}
}
{
    \eR{[English source not present in the checked-in Bencraft/Joly files.]}
}
{
}

\Verse{199}
{
    \zR{妙玉听了这话，想起自家，心上一动，脸 上一热，必然也是红的，倒觉不好意思起来。}
}
{
    \eR{[English source not present in the checked-in Bencraft/Joly files.]}
}
{
}

\Verse{200}
{
    \zR{因站起来说道：“我来得久了，要回 庵里去了。”}
}
{
    \eR{[English source not present in the checked-in Bencraft/Joly files.]}
}
{
}

\Verse{201}
{
    \zR{惜春知妙玉为人，也不深留，送出门口。}
}
{
    \eR{[English source not present in the checked-in Bencraft/Joly files.]}
}
{
}

\Verse{202}
{
    \zR{妙玉笑道：“久已不来，这 里弯弯曲曲的，回去的路头都要迷住了。”}
}
{
    \eR{[English source not present in the checked-in Bencraft/Joly files.]}
}
{
}

\Verse{203}
{
    \zR{宝玉道：“这倒要我来指引指引，何如？”}
}
{
    \eR{[English source not present in the checked-in Bencraft/Joly files.]}
}
{
}

\Verse{204}
{
    \zR{妙玉道：“不敢，二爷前请。”}
}
{
    \eR{[English source not present in the checked-in Bencraft/Joly files.]}
}
{
}

\Verse{205}
{
    \zR{于是二人别了惜春，离了蓼风轩，弯弯曲曲，走近潇湘馆，忽听得叮咚之声。}
}
{
    \eR{[English source not present in the checked-in Bencraft/Joly files.]}
}
{
}

\Verse{206}
{
    \zR{妙玉道：“那里的琴声？”}
}
{
    \eR{[English source not present in the checked-in Bencraft/Joly files.]}
}
{
}

\Verse{207}
{
    \zR{宝玉道：“想必是林妹妹那里抚琴呢。”}
}
{
    \eR{[English source not present in the checked-in Bencraft/Joly files.]}
}
{
}

\Verse{208}
{
    \zR{妙玉道：“原 来他也会这个吗？}
}
{
    \eR{[English source not present in the checked-in Bencraft/Joly files.]}
}
{
}

\Verse{209}
{
    \zR{怎么素日不听见提起？”}
}
{
    \eR{[English source not present in the checked-in Bencraft/Joly files.]}
}
{
}

\Verse{210}
{
    \zR{宝玉悉把黛玉的事说了一遍，因说：“咱 们去看他。”}
}
{
    \eR{[English source not present in the checked-in Bencraft/Joly files.]}
}
{
}

\Verse{211}
{
    \zR{妙玉道：“从古只有听琴，再没有看琴的。”}
}
{
    \eR{[English source not present in the checked-in Bencraft/Joly files.]}
}
{
}

\Verse{212}
{
    \zR{宝玉笑道：“我原说我 是个俗人。”}
}
{
    \eR{[English source not present in the checked-in Bencraft/Joly files.]}
}
{
}

\Verse{213}
{
    \zR{说着，二人走至潇湘馆外，在山子石上坐着静听，甚觉音调清切。}
}
{
    \eR{[English source not present in the checked-in Bencraft/Joly files.]}
}
{
}

\Verse{214}
{
    \zR{只 听得低吟道： 风萧萧兮秋气深，美人千里兮独沉吟。}
}
{
    \eR{[English source not present in the checked-in Bencraft/Joly files.]}
}
{
}

\Verse{215}
{
    \zR{望故乡兮何处？}
}
{
    \eR{[English source not present in the checked-in Bencraft/Joly files.]}
}
{
}

\Verse{216}
{
    \zR{倚栏杆兮涕沾襟。}
}
{
    \eR{[English source not present in the checked-in Bencraft/Joly files.]}
}
{
}

\Verse{217}
{
    \zR{歇了一回，听得又吟道： 山迢迢兮水长，照轩窗兮明月光。}
}
{
    \eR{[English source not present in the checked-in Bencraft/Joly files.]}
}
{
}

\Verse{218}
{
    \zR{耿耿不寐兮银河渺茫，罗衫怯怯兮风露凉。}
}
{
    \eR{[English source not present in the checked-in Bencraft/Joly files.]}
}
{
}

\Verse{219}
{
    \zR{又歇了一歇。}
}
{
    \eR{[English source not present in the checked-in Bencraft/Joly files.]}
}
{
}

\Verse{220}
{
    \zR{妙玉道：“刚才‘侵’字韵是第一叠，如今‘阳’字韵是第二叠了。}
}
{
    \eR{[English source not present in the checked-in Bencraft/Joly files.]}
}
{
}

\Verse{221}
{
    \zR{咱们再听。”}
}
{
    \eR{[English source not present in the checked-in Bencraft/Joly files.]}
}
{
}

\Verse{222}
{
    \zR{里边又吟道： 子之遭兮不自由，予之遇兮多烦忧。}
}
{
    \eR{[English source not present in the checked-in Bencraft/Joly files.]}
}
{
}

\Verse{223}
{
    \zR{之子与我兮心焉相投，思古人兮俾无尤。}
}
{
    \eR{[English source not present in the checked-in Bencraft/Joly files.]}
}
{
}

\Verse{224}
{
    \zR{妙玉道：“这又是一拍。}
}
{
    \eR{[English source not present in the checked-in Bencraft/Joly files.]}
}
{
}

\Verse{225}
{
    \zR{何忧思之深也！”}
}
{
    \eR{[English source not present in the checked-in Bencraft/Joly files.]}
}
{
}

\Verse{226}
{
    \zR{宝玉道：“我虽不懂得，但听他声音， 也觉得过悲了。”}
}
{
    \eR{[English source not present in the checked-in Bencraft/Joly files.]}
}
{
}

\Verse{227}
{
    \zR{里头又调了一回弦。}
}
{
    \eR{[English source not present in the checked-in Bencraft/Joly files.]}
}
{
}

\Verse{228}
{
    \zR{妙玉道：“君弦太高了，与无射律只怕不配 呢。”}
}
{
    \eR{[English source not present in the checked-in Bencraft/Joly files.]}
}
{
}

\Verse{229}
{
    \zR{里边又吟道： 人生斯世兮如轻尘，天上人间兮感夙因。}
}
{
    \eR{[English source not present in the checked-in Bencraft/Joly files.]}
}
{
}

\Verse{230}
{
    \zR{感夙因兮不可？}
}
{
    \eR{[English source not present in the checked-in Bencraft/Joly files.]}
}
{
}

\Verse{231}
{
    \zR{，素心如何天上月! 妙玉听了，呀然失色道：“如何忽作变徵之声？}
}
{
    \eR{[English source not present in the checked-in Bencraft/Joly files.]}
}
{
}

\Verse{232}
{
    \zR{音韵可裂金石矣!只是太过。”}
}
{
    \eR{[English source not present in the checked-in Bencraft/Joly files.]}
}
{
}

\Verse{233}
{
    \zR{宝玉 道：“太过便怎么？”}
}
{
    \eR{[English source not present in the checked-in Bencraft/Joly files.]}
}
{
}

\Verse{234}
{
    \zR{妙玉道：“恐不能持久。”}
}
{
    \eR{[English source not present in the checked-in Bencraft/Joly files.]}
}
{
}

\Verse{235}
{
    \zR{正议论时，听得君弦“蹦”的一 声断了。}
}
{
    \eR{[English source not present in the checked-in Bencraft/Joly files.]}
}
{
}

\Verse{236}
{
    \zR{妙玉站起来，连忙就走。}
}
{
    \eR{[English source not present in the checked-in Bencraft/Joly files.]}
}
{
}

\Verse{237}
{
    \zR{宝玉道：“怎么样？”}
}
{
    \eR{[English source not present in the checked-in Bencraft/Joly files.]}
}
{
}

\Verse{238}
{
    \zR{妙玉道：“日后自知，你 也不必多说。”}
}
{
    \eR{[English source not present in the checked-in Bencraft/Joly files.]}
}
{
}

\Verse{239}
{
    \zR{竟自走了。}
}
{
    \eR{[English source not present in the checked-in Bencraft/Joly files.]}
}
{
}

\Verse{240}
{
    \zR{弄得宝玉满肚疑团，没精打采的，归至怡红院中，不表。}
}
{
    \eR{[English source not present in the checked-in Bencraft/Joly files.]}
}
{
}

\Verse{241}
{
    \zR{且说妙玉归去，早有道婆接着，掩了庵门，坐了一回，把《禅门日诵》念了一 遍。}
}
{
    \eR{[English source not present in the checked-in Bencraft/Joly files.]}
}
{
}

\Verse{242}
{
    \zR{吃了晚饭，点上香，拜了菩萨，命道婆子自去歇着。}
}
{
    \eR{[English source not present in the checked-in Bencraft/Joly files.]}
}
{
}

\Verse{243}
{
    \zR{自己的禅床靠背俱已整齐， 屏息垂帘，跏趺坐下，断除妄想，趁向真如。}
}
{
    \eR{[English source not present in the checked-in Bencraft/Joly files.]}
}
{
}

\Verse{244}
{
    \zR{坐到三更以后，听得房上？？？}
}
{
    \eR{[English source not present in the checked-in Bencraft/Joly files.]}
}
{
}

\Verse{245}
{
    \zR{一片 响声，妙玉恐有贼来，下了禅床，出到前轩，但见云影横空，月华如水。}
}
{
    \eR{[English source not present in the checked-in Bencraft/Joly files.]}
}
{
}

\Verse{246}
{
    \zR{那时天气 尚不很凉，独自一个凭栏站了一回，忽听房上两个猫儿一递一声厮叫。}
}
{
    \eR{[English source not present in the checked-in Bencraft/Joly files.]}
}
{
}

\Verse{247}
{
    \zR{那妙玉忽想 起日间宝玉之言，不觉一阵心跳耳热，自己连忙收摄心神，走进禅房，仍到禅床上 坐了。}
}
{
    \eR{[English source not present in the checked-in Bencraft/Joly files.]}
}
{
}

\Verse{248}
{
    \zR{怎奈神不守舍，一时如万马奔驰，觉得禅床便恍荡起来，身子已不在庵中。}
}
{
    \eR{[English source not present in the checked-in Bencraft/Joly files.]}
}
{
}

\Verse{249}
{
    \zR{便有许多王孙公子，要来娶他；又有些媒婆扯扯拽拽扶他上车，自己不肯去。}
}
{
    \eR{[English source not present in the checked-in Bencraft/Joly files.]}
}
{
}

\Verse{250}
{
    \zR{一回 儿，又有盗贼劫他，持刀执棍的逼勒，只得哭喊求救。}
}
{
    \eR{[English source not present in the checked-in Bencraft/Joly files.]}
}
{
}

\Verse{251}
{
    \zR{早惊醒了庵中女尼道婆等众，都拿火来照看，只见妙玉两手撒开，口中流沫。}
}
{
    \eR{[English source not present in the checked-in Bencraft/Joly files.]}
}
{
}

\Verse{252}
{
    \zR{急叫醒时，只见眼睛直竖，两颧鲜红，骂道：“我是有菩萨保佑，你们这些强徒敢 要怎么样？”}
}
{
    \eR{[English source not present in the checked-in Bencraft/Joly files.]}
}
{
}

\Verse{253}
{
    \zR{众人都唬的没了主意，都说道：“我们在这里呢，快醒转来罢！”}
}
{
    \eR{[English source not present in the checked-in Bencraft/Joly files.]}
}
{
}

\Verse{254}
{
    \zR{妙 玉道：“我要回家去!你们有什么好人，送我回去罢。”}
}
{
    \eR{[English source not present in the checked-in Bencraft/Joly files.]}
}
{
}

\Verse{255}
{
    \zR{道婆道：“这里就是你住 的房子。”}
}
{
    \eR{[English source not present in the checked-in Bencraft/Joly files.]}
}
{
}

\Verse{256}
{
    \zR{说着，又叫别的女尼忙向观音前祷告。}
}
{
    \eR{[English source not present in the checked-in Bencraft/Joly files.]}
}
{
}

\Verse{257}
{
    \zR{求了签，翻开签书看时，是触犯 了西南角上的阴人。}
}
{
    \eR{[English source not present in the checked-in Bencraft/Joly files.]}
}
{
}

\Verse{258}
{
    \zR{就有一个说：“是了，大观园中西南角上本来没有人住，阴气 是有的。”}
}
{
    \eR{[English source not present in the checked-in Bencraft/Joly files.]}
}
{
}

\Verse{259}
{
    \zR{一面弄汤弄水的在那里忙乱。}
}
{
    \eR{[English source not present in the checked-in Bencraft/Joly files.]}
}
{
}

\Verse{260}
{
    \zR{那女尼原是自南边带来的，伏侍妙玉自然 比别人尽心，围着妙玉坐在禅床上。}
}
{
    \eR{[English source not present in the checked-in Bencraft/Joly files.]}
}
{
}

\Verse{261}
{
    \zR{妙玉回头道：“你是谁？”}
}
{
    \eR{[English source not present in the checked-in Bencraft/Joly files.]}
}
{
}

\Verse{262}
{
    \zR{女尼道：“是我。”}
}
{
    \eR{[English source not present in the checked-in Bencraft/Joly files.]}
}
{
}

\Verse{263}
{
    \zR{妙玉仔细瞧了一瞧道：“原来是你！”}
}
{
    \eR{[English source not present in the checked-in Bencraft/Joly files.]}
}
{
}

\Verse{264}
{
    \zR{便抱住那女尼，呜呜咽咽的哭起来，说道： “你是我的妈呀，你不救我，我不得活了！”}
}
{
    \eR{[English source not present in the checked-in Bencraft/Joly files.]}
}
{
}

\Verse{265}
{
    \zR{那女尼一面唤醒他，一面给他揉着。}
}
{
    \eR{[English source not present in the checked-in Bencraft/Joly files.]}
}
{
}

\Verse{266}
{
    \zR{道婆倒上茶来喝了，直到天明才睡了。}
}
{
    \eR{[English source not present in the checked-in Bencraft/Joly files.]}
}
{
}

\Verse{267}
{
    \zR{女尼便打发人去请大夫来看脉。}
}
{
    \eR{[English source not present in the checked-in Bencraft/Joly files.]}
}
{
}

\Verse{268}
{
    \zR{也有说是思虑伤脾的，也有说是热入血室的， 也有说是邪祟触犯的，也有说是内外感冒的：终无定论。}
}
{
    \eR{[English source not present in the checked-in Bencraft/Joly files.]}
}
{
}

\Verse{269}
{
    \zR{后请得一个大夫来看了， 问：“曾打坐过没有？”}
}
{
    \eR{[English source not present in the checked-in Bencraft/Joly files.]}
}
{
}

\Verse{270}
{
    \zR{道婆说道：“向来打坐的。”}
}
{
    \eR{[English source not present in the checked-in Bencraft/Joly files.]}
}
{
}

\Verse{271}
{
    \zR{大夫道：“这病可是昨夜忽 然来的么？”}
}
{
    \eR{[English source not present in the checked-in Bencraft/Joly files.]}
}
{
}

\Verse{272}
{
    \zR{道婆道：“是。”}
}
{
    \eR{[English source not present in the checked-in Bencraft/Joly files.]}
}
{
}

\Verse{273}
{
    \zR{大夫道：“这是走魔入火的原故。”}
}
{
    \eR{[English source not present in the checked-in Bencraft/Joly files.]}
}
{
}

\Verse{274}
{
    \zR{众人问：“有 碍没有？”}
}
{
    \eR{[English source not present in the checked-in Bencraft/Joly files.]}
}
{
}

\Verse{275}
{
    \zR{大夫道：“幸亏打坐不久，魔还入得浅，可以有救。”}
}
{
    \eR{[English source not present in the checked-in Bencraft/Joly files.]}
}
{
}

\Verse{276}
{
    \zR{写了降伏心火的 药，吃了一剂，稍稍平复些。}
}
{
    \eR{[English source not present in the checked-in Bencraft/Joly files.]}
}
{
}

\Verse{277}
{
    \zR{外面那些游头浪子听见了，便造作许多谣言，说：“这 么年纪，那里忍得住？}
}
{
    \eR{[English source not present in the checked-in Bencraft/Joly files.]}
}
{
}

\Verse{278}
{
    \zR{况且又是很风流的人品，很乖觉的性灵!以后不知飞在谁手 里，便宜谁去呢。”}
}
{
    \eR{[English source not present in the checked-in Bencraft/Joly files.]}
}
{
}

\Verse{279}
{
    \zR{过了几日，妙玉病虽略好了些，神思未复，终有些恍惚。}
}
{
    \eR{[English source not present in the checked-in Bencraft/Joly files.]}
}
{
}

\Verse{280}
{
    \zR{一日，惜春正坐着，彩屏忽然进来，回道：“姑娘知道妙玉师父的事吗？”}
}
{
    \eR{[English source not present in the checked-in Bencraft/Joly files.]}
}
{
}

\Verse{281}
{
    \zR{惜 春道：“他有什么事？”}
}
{
    \eR{[English source not present in the checked-in Bencraft/Joly files.]}
}
{
}

\Verse{282}
{
    \zR{彩屏道：“我昨日听见邢姑娘和大奶奶在那里说呢：他自 从那日合姑娘下棋回去，夜间忽然中了邪，嘴里乱嚷，说强盗来抢他来了。}
}
{
    \eR{[English source not present in the checked-in Bencraft/Joly files.]}
}
{
}

\Verse{283}
{
    \zR{到如今 还没好呢。}
}
{
    \eR{[English source not present in the checked-in Bencraft/Joly files.]}
}
{
}

\Verse{284}
{
    \zR{姑娘，你说这不是奇事吗？”}
}
{
    \eR{[English source not present in the checked-in Bencraft/Joly files.]}
}
{
}

\Verse{285}
{
    \zR{惜春听了，默默无语。}
}
{
    \eR{[English source not present in the checked-in Bencraft/Joly files.]}
}
{
}

\Verse{286}
{
    \zR{因想：“妙玉虽然 洁净，毕竟尘缘未断。}
}
{
    \eR{[English source not present in the checked-in Bencraft/Joly files.]}
}
{
}

\Verse{287}
{
    \zR{可惜我生在这种人家，不便出家，我若出了家时，那有邪魔 缠扰？}
}
{
    \eR{[English source not present in the checked-in Bencraft/Joly files.]}
}
{
}

\Verse{288}
{
    \zR{一念不生，万缘俱寂。”}
}
{
    \eR{[English source not present in the checked-in Bencraft/Joly files.]}
}
{
}

\Verse{289}
{
    \zR{想到这里，蓦与神会，若有所得，便口占一偈云： 大造本无方，云何是应住？}
}
{
    \eR{[English source not present in the checked-in Bencraft/Joly files.]}
}
{
}

\Verse{290}
{
    \zR{既从空中来，应向空中去。}
}
{
    \eR{[English source not present in the checked-in Bencraft/Joly files.]}
}
{
}

\Verse{291}
{
    \zR{占毕，即命丫头焚香。}
}
{
    \eR{[English source not present in the checked-in Bencraft/Joly files.]}
}
{
}

\Verse{292}
{
    \zR{自己静坐了一回，又翻开那棋谱来，把孔融、王积薪等 所著看了几篇。}
}
{
    \eR{[English source not present in the checked-in Bencraft/Joly files.]}
}
{
}

\Verse{293}
{
    \zR{内中“茂叶包蟹势”、“黄莺搏兔势”，都不出奇；“三十六局杀 角势”，一时也难会难记；独看到“十龙走马”，觉得甚有意思。}
}
{
    \eR{[English source not present in the checked-in Bencraft/Joly files.]}
}
{
}

\Verse{294}
{
    \zR{正在那里作想， 只听见外面一个人走进院来，连叫彩屏。}
}
{
    \eR{[English source not present in the checked-in Bencraft/Joly files.]}
}
{
}

\Verse{295}
{
    \zR{未知是谁，下回分解。}
}
{
    \eR{[English source not present in the checked-in Bencraft/Joly files.]}
}
{
}

\Verse{296}
{
    \zR{(本章完)}
}
{
    \eR{[English source not present in the checked-in Bencraft/Joly files.]}
}
{
}
